% Zone „Das kennst du schon“ – aus bank/konfidenzintervalle/zone.jsonl (zusammenbau v0.1)
\einheitenkopf[zone][Kennst du schon]{Das kennst du schon}
\setcounter{aufgabe}{0}

%% TODO Titel: Ich-kann-Satz fehlt in der Bank (Platzhalter: Kettenname) – Nr. 1
%% TODO Anweisung: ein Satz über der ersten Teilaufgabe fehlt in der Bank – Nr. 1
\begin{aufgabe}{Relative Häufigkeit und Anteil aus einer Stichprobe bilden}
\begin{teile}
\teil 36 von 120 Befragten – welcher Anteil? \leerfeld
\teil Von 1250 Befragten lesen 287 täglich Zeitung. Gib den Anteil auf Tausendstel gerundet an. \leerfeld
\end{teile}
\end{aufgabe}

%% TODO Titel: Ich-kann-Satz fehlt in der Bank (Platzhalter: Kettenname) – Nr. 2
%% TODO Anweisung: ein Satz über der ersten Teilaufgabe fehlt in der Bank – Nr. 2
\begin{aufgabe}{Binomialverteilung ansetzen und Punktwahrscheinlichkeiten berechnen}
\begin{teile}
\teil Eine Münze wird dreimal geworfen, X zählt „Kopf“. Berechne $P(X = 3)$. \leerfeld
\teil X ist binomialverteilt mit $n = 10$ und $p = 0{,}8$. Berechne $P(X = 9)$ auf Tausendstel. \leerfeld
\end{teile}
\end{aufgabe}

%% TODO Titel: Ich-kann-Satz fehlt in der Bank (Platzhalter: Kettenname) – Nr. 3
%% TODO Anweisung: ein Satz über der ersten Teilaufgabe fehlt in der Bank – Nr. 3
\begin{aufgabe}{Graphen zweier Funktionen ablesen (Schnitt mit Waagerechten und Senkrechten)}
\begin{teile}
\teil Lies ab: An welcher Stelle x hat f den Wert 3?

\begin{ksys}[xmin=0,xmax=5,ymin=0,ymax=5,ablesen] \gerade{1}{1}{f} \end{ksys}
x = \leerfeld
\teil Lies ab: An welchen Stellen x haben f und g jeweils den Wert 3?

\begin{ksys}[xmin=0,xmax=5,ymin=0,ymax=5,ablesen] \gerade{2}{-1}{f} \gerade{0.5}{1}{g} \end{ksys}
f: x = \leerfeld , g: x = \leerfeld
\end{teile}
\end{aufgabe}

%% TODO Titel: Ich-kann-Satz fehlt in der Bank (Platzhalter: Kettenname) – Nr. 4
%% TODO Anweisung: ein Satz über der ersten Teilaufgabe fehlt in der Bank – Nr. 4
\begin{aufgabe}{Die Sigma-Regeln und ihre c-Werte}
\begin{teile}
\teil Wie viel Prozent der Werte einer Binomialverteilung liegen nach den Sigma-Regeln etwa im Intervall von $\mu - 1{,}96\sigma$ bis $\mu + 1{,}96\sigma$? \leerfeld \%
\teil X ist binomialverteilt mit $n = 100$ und $p = 0{,}5$. Berechne $\mu$, $\sigma$ und die Grenzen $\mu \pm 1{,}96\sigma$.
\end{teile}
\end{aufgabe}

%% TODO Titel: Ich-kann-Satz fehlt in der Bank (Platzhalter: Kettenname) – Nr. 5
%% TODO Anweisung: ein Satz über der ersten Teilaufgabe fehlt in der Bank – Nr. 5
\begin{aufgabe}{Quadratische Gleichungen und Wurzelgleichungen lösen, Lösungen deuten}
\begin{gleichungsraster}[2]
\gl{\text{Löse $x^2 = 0{,}09$.}} &
\gl{\text{Löse $x^2 - 0{,}8x + 0{,}15 = 0$.}} \\
\end{gleichungsraster}
\end{aufgabe}

%% TODO Titel: Ich-kann-Satz fehlt in der Bank (Platzhalter: Kettenname) – Nr. 6
%% TODO Anweisung: ein Satz über der ersten Teilaufgabe fehlt in der Bank – Nr. 6
\begin{aufgabe}{Ungleichungen in n lösen und runden}
\begin{teile}
\teil Für welche kleinste natürliche Zahl n gilt $\sqrt{n} > 10$? n = \leerfeld
\teil Für welche kleinste natürliche Zahl n gilt $\frac{0{,}4}{\sqrt{n}} \le 0{,}02$? n = \leerfeld
\end{teile}
\end{aufgabe}

%% TODO Titel: Ich-kann-Satz fehlt in der Bank (Platzhalter: Kettenname) – Nr. 7
%% TODO Anweisung: ein Satz über der ersten Teilaufgabe fehlt in der Bank – Nr. 7
\begin{aufgabe}{Binomialverteilung ansetzen und Punktwahrscheinlichkeiten berechnen – Fehler finden}
\begin{teile}
\teil X ist binomialverteilt mit $n = 10$ und $p = 0{,}9$. Ali soll $P(X = 9)$ berechnen und rechnet $P(X = 9) + P(X = 10) \approx 0{,}736$. Finde den Fehler und rechne richtig.
\end{teile}
\end{aufgabe}

%% TODO Titel: Ich-kann-Satz fehlt in der Bank (Platzhalter: Kettenname) – Nr. 8
%% TODO Anweisung: ein Satz über der ersten Teilaufgabe fehlt in der Bank – Nr. 8
\begin{aufgabe}{Binomialverteilung ansetzen und Punktwahrscheinlichkeiten berechnen – selbst rechnen}
\begin{teile}
\teil X ist binomialverteilt mit $n = 12$ und $p = 0{,}9$. Berechne $P(X = 11)$ auf Tausendstel. \leerfeld
\end{teile}
\end{aufgabe}
